%% LyX 2.4.4 created this file.  For more info, see https://www.lyx.org/.
%% Do not edit unless you really know what you are doing.
\documentclass[english]{article}
\usepackage[T1]{fontenc}
\usepackage[latin9]{inputenc}
\usepackage{enumitem}
\usepackage{amsmath}
\usepackage{amsthm}
\usepackage{cancel}
\usepackage{geometry}
\geometry{verbose,tmargin=1in,bmargin=1in,lmargin=1in,rmargin=1in}

\makeatletter
%%%%%%%%%%%%%%%%%%%%%%%%%%%%%% Textclass specific LaTeX commands.
\numberwithin{equation}{section}
\numberwithin{figure}{section}
\newlength{\lyxlabelwidth}      % auxiliary length 

%%%%%%%%%%%%%%%%%%%%%%%%%%%%%% User specified LaTeX commands.
\usepackage{fancyhdr}
\usepackage{lastpage}
\pagestyle{fancy}
\usepackage{enumitem}
\usepackage{ifthen}
%sets math fonts to be more readable
\everymath=\expandafter{\the\everymath\displaystyle}
%\DeclareMathSizes{10}{10}{10}{10}
\usepackage{multicol}

\usepackage{tikz}
\usetikzlibrary{plotmarks}
\usepackage{pgfplots}
\pgfplotsset{compat=newest}

\usepackage{calc}
\usepackage[nomessages]{fp}

\usepackage{advdate}    % Advancing/saving dates
\usepackage{datetime}   % Dates formatting
\usepackage{datenumber} % Counters for dates
% Added by lyx2lyx
\setlength{\parskip}{\smallskipamount}
\setlength{\parindent}{0pt}

\makeatother

\usepackage{babel}
\begin{document}
\renewcommand{\author}{Dr.\,James Ashe}

\newcommand{\classNum}{131}

\newcommand{\classSec}{B}

\newcommand{\nameType}{Name:}

\newcommand{\examType}{Exam 4} 

\renewcommand{\date}{\formatdate{18}{11}{2013}}

%%First page's header%%%%%%%%%%%%%%%%%%%%%%
\fancypagestyle{firststyle} %%header settings for first pagr
{
	\fancyhead[L]{MTH \classNum{}(\classSec{}) \author{}\\
	\textbf{\examType{}} \date{}}
	\fancyhead[C]{\textbf{\nameType}}
	\setlength{\headsep}{14pt}
}
%%%%%%%%%%%%%%%%%%%%%%%%%%%%%%%%%%%%%%%%%%

%%Next page's header%%%%%%%%%%%%%%%%%%%%%%%
\setlist{nolistsep}
\lhead{MTH \classNum{}(\classSec{})} 
\chead{\textbf{\nameType}} 
\cfoot{\thepage{}~of~\pageref{LastPage}} 
\setlength{\headsep}{5pt} 
%%%%%%%%%%%%%%%%%%%%%%%%%%%%%%%%%%%%%%%%%%%

%%Various settings; see the preamble for other settings%%%%%
%\setenumerate[0]{leftmargin=12pt,itemindent=0pt,labelwidth=0pt}
%\setlist[enumerate]{topsep=0pt,itemsep=1pt,partopsep=1ex,parsep=1ex}
\renewcommand{\labelenumi}{\textbf{\arabic{enumi}.}}
\renewcommand{\labelenumii}{\textbf{\alph{enumii}.}}
\thispagestyle{firststyle}
%%%%%%%%%%%%%%%%%%%%%%%%%%%%%%%%%%%%%%%%%%
\newcommand{\xmax}{10}
\newcommand{\xmin}{-10}
\newcommand{\xstep}{0} %must be xmin+n
\newcommand{\ymax}{10}
\newcommand{\ymin}{-10}
\newcommand{\ystep}{0} %must be ymin+n

\newcommand{\x}{2} \newcommand{\y}{1} \newcommand{\z}{1} 
\newcommand{\Ro}{1} \newcommand{\Po}{1} \newcommand{\Ao}{1} \newcommand{\To}{1}

\newcounter{probs}
\newcommand{\xspace}{\vspace{1.5cm}}

\textbf{You must show your work to recieve. Each problem is worth
4 pts unless stated otherwise.}
\begin{enumerate}[resume]
\item \renewcommand{\x}{2} 
\renewcommand{\y}{1} 
\renewcommand{\z}{1}Use the function, $f(x)=\x^{x}$, to complete parts \ref{enu:Graph-start}--\ref{enu:Find-end}.\begin{multicols}{2}

\begin{enumerate}
\item \label{enu:Graph-start}Graph $f(x)$. Label the \emph{y}-int and
at least one other point.\addtocounter{probs}{1}
\item Find the domain of $f(x)$.

\begin{enumerate}
\item []$(-\infty,\infty)$\vfill
\end{enumerate}
\item Find the range of $f(x)$.

\begin{enumerate}
\item []$(0,\infty)$\vfill
\end{enumerate}
\item \label{enu:Find-end}Find the horizontal asymptote of $f(x)$.

\begin{enumerate}
\item []$y=0$\vfill\columnbreak
\end{enumerate}
\item []\renewcommand{\xmax}{10} 
\renewcommand{\xstep}{5}
\renewcommand{\ymax}{\xmax} 
\renewcommand{\ystep}{5}
%%%%%%%
\begin{tikzpicture} 
	\begin{axis}[
		axis x line=middle, axis y line=middle,     		
		xtick={-\xmax,-\xstep,...,\xmax},
		ytick={-\ymax,-\ystep,...,\ymax},
		minor x tick num={4},
		minor y tick num={4},
		grid=both,
		xlabel={$$},     
		ylabel={$$},
		xmin=-\xmax.25,xmax=\xmax.25,
		ymin=-\ymax.25,ymax=\ymax.25,     		width=5.5cm,     		height=5.5cm] 		
\addplot[red,very thick,mark=noner,smooth,domain=\xmin:\xmax]{2^x}; 
	\end{axis}
\end{tikzpicture}
\end{enumerate}
\end{multicols}
\item \renewcommand{\x}{2} 
\renewcommand{\y}{1} 
\renewcommand{\z}{5}Use the function, $f(x)=\left(\frac{1}{\x}\right)^{x}$, to complete
parts \ref{enu:Graph-start-1-1}--\ref{enu:Find-end-1-1}.\begin{multicols}{2}

\begin{enumerate}
\item \label{enu:Graph-start-1-1}Graph $f(x)$. Label the \emph{y}-int
and at least one other point.\addtocounter{probs}{1}
\item \label{enu:Find-end-1-1}Find the horizontal asymptote of \\
$g(x)=\left(\frac{1}{\x}\right)^{x}-\z$.

\begin{enumerate}
\item []$y=0$\vfill\addtocounter{probs}{1}\columnbreak
\end{enumerate}
\item []\renewcommand{\xmax}{10} 
\renewcommand{\xstep}{5}
\renewcommand{\ymax}{\xmax} 
\renewcommand{\ystep}{5}
%%%%%%%
\begin{tikzpicture} 
	\begin{axis}[
		axis x line=middle, axis y line=middle,     		
		xtick={-\xmax,-\xstep,...,\xmax},
		ytick={-\ymax,-\ystep,...,\ymax},
		minor x tick num={4},
		minor y tick num={4},
		grid=both,
		xlabel={$$},     
		ylabel={$$},
		xmin=-\xmax.25,xmax=\xmax.25,
		ymin=-\ymax.25,ymax=\ymax.25,     		width=5.5cm,     		height=5.5cm] 		
\addplot[red,very thick,mark=noner,smooth,domain=\xmin:\xmax]{(1/2)^x}; 
	\end{axis}
\end{tikzpicture}
\end{enumerate}
\end{multicols}
\item \renewcommand{\x}{2} 
\renewcommand{\y}{4} 
\renewcommand{\z}{1}Use the graph of $b^{x}$ in parts \ref{enu:Graph-start-1}--\ref{enu:Find-end-1}
to graph the given transformation.\begin{multicols}{3} \newcommand{\xWidth}{6cm}

\begin{enumerate}
\item \label{enu:Graph-start-1}Graph $b^{x}+\x$\addtocounter{probs}{1}
\item [] \hspace{-1.5cm}\renewcommand{\xmax}{10} 
\renewcommand{\xstep}{5}
\renewcommand{\ymax}{\xmax} 
\renewcommand{\ystep}{5}
%%%%%%%
\begin{tikzpicture} 
	\begin{axis}[
		axis x line=middle, axis y line=middle,     		
		xtick={-\xmax,-\xstep,...,\xmax},
		ytick={-\ymax,-\ystep,...,\ymax},
		minor x tick num={4},
		minor y tick num={4},
		grid=both,
		xlabel={$$},     
		ylabel={$$},
		xmin=-\xmax.25,xmax=\xmax.25,
		ymin=-\ymax.25,ymax=\ymax.25,     		width=\xWidth,     		height=\xWidth] 		
\addplot[red,very thick,mark=noner,smooth,domain=\xmin:\xmax]{2^x};
\addplot[blue,dashed,very thick,mark=noner,smooth,domain=\xmin:\xmax]{2^x+2};
	\end{axis}
\end{tikzpicture}\columnbreak
\item Graph $b^{x+\y}$\addtocounter{probs}{1}
\item [] \hspace{-1.5cm}\renewcommand{\xmax}{10} 
\renewcommand{\xstep}{5}
\renewcommand{\ymax}{\xmax} 
\renewcommand{\ystep}{5}
%%%%%%%
\begin{tikzpicture} 
	\begin{axis}[
		axis x line=middle, axis y line=middle,     		
		xtick={-\xmax,-\xstep,...,\xmax},
		ytick={-\ymax,-\ystep,...,\ymax},
		minor x tick num={4},
		minor y tick num={4},
		grid=both,
		xlabel={$$},     
		ylabel={$$},
		xmin=-\xmax.25,xmax=\xmax.25,
		ymin=-\ymax.25,ymax=\ymax.25,     		width=\xWidth,     		height=\xWidth] 		
\addplot[red,very thick,mark=noner,smooth,domain=\xmin:\xmax]{2^x};
\addplot[blue,dashed,very thick,mark=noner,smooth,domain=\xmin:\xmax]{2^x};
	\end{axis}
\end{tikzpicture}\columnbreak
\item \label{enu:Find-end-1}Graph $b^{x+\y}+\x$\addtocounter{probs}{1}
\item [] \hspace{-1.5cm}\renewcommand{\xmax}{10} 
\renewcommand{\xstep}{5}
\renewcommand{\ymax}{\xmax} 
\renewcommand{\ystep}{5}
%%%%%%%
\begin{tikzpicture} 
	\begin{axis}[
		axis x line=middle, axis y line=middle,     		
		xtick={-\xmax,-\xstep,...,\xmax},
		ytick={-\ymax,-\ystep,...,\ymax},
		minor x tick num={4},
		minor y tick num={4},
		grid=both,
		xlabel={$$},     
		ylabel={$$},
		xmin=-\xmax.25,xmax=\xmax.25,
		ymin=-\ymax.25,ymax=\ymax.25,     		width=\xWidth,     		height=\xWidth] 		
\addplot[red,very thick,mark=noner,smooth,domain=\xmin:\xmax]{2^x};
%\addplot[red,very thick,mark=noner,smooth,domain=\xmin:\xmax]{2^x};
	\end{axis}
\end{tikzpicture}\columnbreak
\end{enumerate}
\end{multicols}
\end{enumerate}
\emph{Use the formulas, $A=P(1+r)^{t}$, $A=P(1+\frac{r}{n})^{nt}$,
or $A=Pe^{rt}$ to problems \ref{enu:first}--\ref{enu:last}.}
\begin{enumerate}[resume]
\item \label{enu:first}{[}6 points{]}\renewcommand{\Po}{850} 
\renewcommand{\Ro}{9} 
\renewcommand{\To}{10}
\renewcommand{\Ao}{1} Find the accumulated value of an investment of \$\Po~at \Ro\% interest
compounded quarterly after \To~ years.\vfill\addtocounter{probs}{1}
\end{enumerate}
\newpage
\begin{enumerate}[resume]
\item \renewcommand{\Po}{150,000} 
\renewcommand{\Ro}{4.5} 
\renewcommand{\To}{15}
\renewcommand{\Ao}{10} {[}6 points{]} Bob buys a house for \$\Po. The value of the house
grows at a rate of \Ro\% compounded monthly. Find how much the house
is worth after \To~years.\vspace{1cm}\vfill\addtocounter{probs}{1}
\item \label{enu:last}\renewcommand{\Po}{2,000} 
\renewcommand{\Ro}{6.5} 
\renewcommand{\To}{10}
\renewcommand{\Ao}{1} {[}6 points{]} Suppose Jane has \$\Po to invest. Which investment
yields the greater return over \To years: 

\begin{enumerate}[resume]
\item []7\% compounded monthly or \Ro\% compounded continuously? Justiy
your answer.\vfill\vspace{2cm}\addtocounter{probs}{1}
\end{enumerate}
\item Write the expression in its equivalent \emph{exponential} form for
parts \ref{enu:first-1}--\ref{enu:last-1}.\begin{multicols}{3}

\begin{enumerate}
\item \label{enu:first-1}\renewcommand{\x}{x} 
\renewcommand{\y}{9} 
\renewcommand{\z}{3}$\log_{\z}\x=\y$\xspace\addtocounter{probs}{1}
\item \renewcommand{\x}{5} 
\renewcommand{\y}{1} 
\renewcommand{\z}{1}$\ln x=\x$\xspace\addtocounter{probs}{1}
\item \label{enu:last-1}\renewcommand{\x}{2} 
\renewcommand{\y}{c} 
\renewcommand{\z}{x}$\log_{\z}\x=\y$\xspace\addtocounter{probs}{1}
\end{enumerate}
\end{multicols}\xspace
\item Write the expression in its equivalent \emph{logarithmic} form for
parts \ref{enu:first-1-1}--\ref{enu:last-1-1}.\begin{multicols}{3}

\begin{enumerate}
\item \label{enu:first-1-1}\renewcommand{\x}{x} 
\renewcommand{\y}{125} 
\renewcommand{\z}{5}$\z^{\x}=\y$\xspace\addtocounter{probs}{1}
\item \renewcommand{\x}{12} 
\renewcommand{\y}{7} 
\renewcommand{\z}{x}$\z^{\x}=\y$\xspace\addtocounter{probs}{1}
\item \label{enu:last-1-1}\renewcommand{\x}{-2} 
\renewcommand{\y}{\frac{1}{16}} 
\renewcommand{\z}{x}$\z^{\x}=\y$\xspace\addtocounter{probs}{1}
\end{enumerate}
\end{multicols}\xspace
\item {[}3 points each{]} Evaluate or simplify the expression without using
a calculator for parts \ref{enu:first-1-2}--\ref{enu:last-1-2}.\begin{multicols}{3}

\begin{enumerate}
\item \label{enu:first-1-2}\renewcommand{\x}{2} 
\renewcommand{\y}{1} 
\renewcommand{\z}{1}$\log_{5}5$\xspace\addtocounter{probs}{1}
\item \renewcommand{\x}{2} 
\renewcommand{\y}{1} 
\renewcommand{\z}{1}$e^{\ln2}$\xspace\addtocounter{probs}{1}
\item \renewcommand{\x}{2} 
\renewcommand{\y}{1} 
\renewcommand{\z}{1}$\log_{2}2^{10}$\xspace\addtocounter{probs}{1}
\item \renewcommand{\x}{2} 
\renewcommand{\y}{1} 
\renewcommand{\z}{1}$4^{\log_{4}11}$\xspace\addtocounter{probs}{1}
\item \renewcommand{\x}{2} 
\renewcommand{\y}{1} 
\renewcommand{\z}{1}$\log_{5}25$\xspace\addtocounter{probs}{1}
\item \label{enu:last-1-2}\renewcommand{\x}{2} 
\renewcommand{\y}{1} 
\renewcommand{\z}{1}$\log10$\xspace\vfill
\end{enumerate}
\end{multicols}
\item Find the domain of the logarithmic function for parts \ref{enu:first-1-2-1}--\ref{enu:last-1-2-1}.\begin{multicols}{2}

\begin{enumerate}
\item \label{enu:first-1-2-1}\renewcommand{\x}{2} 
\renewcommand{\y}{1} 
\renewcommand{\z}{1}$\log_{13}x$\addtocounter{probs}{1}
\item \label{enu:last-1-2-1}\renewcommand{\x}{2} 
\renewcommand{\y}{1} 
\renewcommand{\z}{1}$\ln(x-2)$\addtocounter{probs}{1}
\end{enumerate}
\end{multicols}\xspace
\end{enumerate}

\end{document}
